% Part: second-order-logic
% Chapter: sol-and-sets
% Section: power-of-continuum

\documentclass[../../../include/open-logic-section]{subfiles}

\begin{document}

\olfileid[hi]{sol}{set}{pow}

\olsection{सातत्य की शक्ति}

\begin{explain}
द्वितीय-कोटि तर्कशास्त्र में हम !!{domain} के उपसमुच्चयों पर परिमाणीकरण कर
सकते हैं, पर !!{domain} के उपसमुच्चयों के समुच्चयों पर नहीं। सीधे ऐसा करने
के लिए हमें \emph{तृतीय-कोटि} तर्कशास्त्र चाहिए होगा। उदाहरणार्थ, यदि हम
कैंटर का यह प्रमेय कहना चाहें कि किसी समुच्चय के घात-समुच्चय से उसी समुच्चय
में कोई !!{injective} फलन नहीं होता, तो हम इसे इस प्रकार निरूपित करने का
प्रयास कर सकते हैं: ``प्रत्येक समुच्चय~$X$ और प्रत्येक समुच्चय~$P$ के लिए,
यदि $P$, ~$X$ का घात-समुच्चय है, तो $\cardle{P}{X}$ नहीं।'' और यह कहने के
लिए कि $P$, ~$X$ का घात-समुच्चय है, हमें औपचारिक रूप से व्यक्त करना होगा
कि $P$ के !!{element} ठीक वही और केवल वही हैं जो~$X$ के उपसमुच्चय हैं;
अर्थात् $\lforall[Y][(P(Y) \liff Y \subseteq X)]$ जैसा कुछ। समस्या
$P(Y)$ में है: यह द्वितीय-कोटि तर्कशास्त्र का !!a{formula} नहीं है, क्योंकि
~$P$ जैसे एक-स्थानी संबंध-चरों के तर्क केवल पद हो सकते हैं।

फिर भी, यदि प्रांत पर्याप्त बड़ा हो, तो हम समुच्चयों के समुच्चयों पर
परिमाणीकरण का \emph{अनुकरण} कर सकते हैं। विचार इस तथ्य का उपयोग करना है कि
दो-स्थानी संबंध~$R$, !!{domain} के !!{element} को उसी !!{domain} के
!!{element} से संबद्ध करता है। ऐसा~$R$ दिए होने पर हम उन सभी !!{element}
को एकत्र कर सकते हैं जिनसे कोई~$x$, ~$R$-संबद्ध है: $\Setabs{y \in
  \Domain{M}}{R(x,y)}$ वह समुच्चय है जिसे~$x$ ``कूटित करता है''। उलटे, यदि
$Z \subseteq \Pow{\Domain{M}}$, ~$\Domain{M}$ के उपसमुच्चयों का कोई संग्रह
है, और~$\Domain{M}$ में कम-से-कम उतने !!{element} हैं जितने~$Z$ में
समुच्चय, तो कोई ऐसा संबंध~$R \subseteq \Domain{M}^2$ भी है कि~$Z$ का
प्रत्येक $Y$ किसी~$x$ द्वारा~$R$ के माध्यम से कूटित होता है।
\end{explain}

\begin{defn}
यदि $R \subseteq \Domain{M}^2$, तो \emph{$x$, ~$R$ के द्वारा}
$\Setabs{y \in \Domain{M}}{R(x, y)}$ \emph{को कूटित करता है}।
\end{defn}

यदि कोई !!a{element}~$x \in \Domain{M}$, ~$R$ के द्वारा किसी समुच्चय $Z
\subseteq \Domain{M}$ को कूटित करता है, तो समुच्चय $Y \subseteq
\Domain{M}$ समुच्चयों के एक समुच्चय को कूटित करता है---अर्थात् उन
समुच्चयों को जिन्हें~$Y$ के !!{element} कूटित करते हैं। इसलिए कोई समुच्चय
$Y$, ~$\Pow{X}$ को~$R$ द्वारा कूटित कर सकता है। वह ऐसा तभी और केवल तभी
करता है जब प्रत्येक $Z \subseteq X$ को कोई $x \in Y$, ~$R$ द्वारा कूटित
करता हो और प्रत्येक $x \in Y$, किसी $Z \subseteq X$ को~$R$ द्वारा कूटित
करता हो।

\begin{prop}
!!{formula}
  \[
  \fn{Codes}(x, R, Z) \ident \lforall[y][(Z(y) \liff
    R(x, y))]
  \]
यह व्यक्त करता है कि $s(x)$, ~$s(R)$ के द्वारा $s(Z)$ को कूटित करता है।
!!{formula}
\begin{multline*}
  \fn{Pow}(Y, R, X) \ident \\
  \lforall[Z][(Z \subseteq X \lif \lexists[x][(Y(x) \land
    \fn{Codes}(x, R, Z))])] \land {} \\ \lforall[x][(Y(x) \lif
  \lforall[Z][(\fn{Codes}(x, R, Z) \lif Z \subseteq X)]]
\end{multline*}
यह व्यक्त करता है कि $s(Y)$, ~$s(R)$ के द्वारा~$s(X)$ के घात-समुच्चय को
कूटित करता है; अर्थात् $s(Y)$ के अवयव~$s(R)$ द्वारा ठीक~$s(X)$ के
उपसमुच्चयों को कूटित करते हैं।
\end{prop}

\begin{explain}
इस युक्ति से हम स्वयं उपसमुच्चयों के बजाय उनके कूटों पर परिमाणीकरण करके
घात-समुच्चय के विषय में कथन व्यक्त कर सकते हैं। उदाहरणार्थ, अब कैंटर के
प्रमेय को यह कहकर व्यक्त किया जा सकता है कि~$X$ के घात-समुच्चय को कूटित
करने वाले किसी भी संबंध के प्रांत से स्वयं~$X$ में कोई !!{injective} फलन
नहीं है।
\end{explain}

\begin{prop}
यह वाक्य
\begin{multline*}
  \lforall[X][\lforall[Y][\lforall[R][(\fn{Pow}(Y, R, X) \lif \\
      \lnot \lexists[u][(\lforall[x][\lforall[y][(\eq[u(x)][u(y)] \lif
            \eq[x][y])]] \land {}\\
        \lforall[x][(Y(x) \lif X(u(x)))])])]]]
\end{multline*}
वैध है।
\end{prop}

\begin{explain}
किसी !!a{denumerable} समुच्चय का घात-समुच्चय !!{nonenumerable} होता है,
इसलिए उसकी कार्डिनैलिटी किसी भी !!{denumerable} समुच्चय की कार्डिनैलिटी
(जो $\aleph_0$ है) से बड़ी होती है। $\Pow{\Nat}$ के आकार को ``सातत्य की
शक्ति'' कहते हैं, क्योंकि वह वास्तविक संख्या-रेखा~$\Real$ के बिंदुओं के
समुच्चय के समान आकार का है। यदि !!{domain} इतना बड़ा है कि किसी
!!a{denumerable} समुच्चय के घात-समुच्चय को कूटित कर सके, तो हम यह बात इस
प्रकार व्यक्त कर सकते हैं कि कोई समुच्चय सातत्य के आकार का है: वह ऐसे किसी
समुच्चय~$Y$ के समसंख्यक है जो आकार~$\aleph_0$ वाले समुच्चय~$X$ के
घात-समुच्चय को कूटित करता है। (यदि प्रांत पर्याप्त बड़ा नहीं है, अर्थात्
उसमें~$\Real$ के समसंख्यक कोई उपसमुच्चय नहीं है, तो~$\Pow{X}$ को कूटित करने
वाला कोई संबंध भी नहीं हो सकता।)
\end{explain}

\begin{prop}
यदि $\cardle{\Real}{\Domain{M}}$, तो !!{formula}
\begin{multline*}
\fn{Cont}(Y) \ident 
\lexists[X][\lexists[R][((\fn{Aleph}_0(X) \land
      \fn{Pow}(Y, R, X)) \land \\ 
      \lforall[x][\lforall[y][((Y(x) \land {} 
      Y(y) \land \lforall[z][R(x,z) \liff R(y,z)]) \lif \eq[x][y])]])]]
\end{multline*}
यह व्यक्त करता है कि $\cardeq{s(Y)}{\Real}$।
\end{prop}

\begin{proof}
  $\fn{Pow}(Y, R, X)$ व्यक्त करता है कि $s(Y)$, ~$s(R)$ के द्वारा~$s(X)$
  के घात-समुच्चय को कूटित करता है, और $\fn{Aleph}_0(X)$ कहता है कि~$s(X)$
  परिगणनीय है। अतः $s(Y)$ कम-से-कम सातत्य की शक्ति जितना बड़ा है, यद्यपि
  वह उससे बड़ा भी हो सकता है (यदि $s(Y)$ के अनेक अवयव~$X$ के एक ही
  उपसमुच्चय को कूटित करते हों)। अंतिम संयोज्य इसे निषिद्ध करता है; वह माँग
  करता है कि $s(R)$ के माध्यम से $s(Y)$ के अवयवों और $s(Z)$ के उपसमुच्चयों
  के बीच संबंध !!{injective} हो।
\end{proof}

\begin{prop}
$\cardeq{\Domain{M}}{\Real}$ तभी और केवल तभी जब
\begin{multline*}
  \Sat{M}{\lexists[X][\lexists[Y][\lexists[R][(\fn{Aleph}_0(X) \land \fn{Pow}(Y, R, X) \land \\
          \lexists[u][(\lforall[x][\lforall[y][(\eq[u(x)][u(y)] \lif \eq[x][y])]] \land {} \\\lforall[y][(Y(y) \lif \lexists[x][\eq[y][u(x)]])])])]]]}.
        \end{multline*}
  \end{prop}

\begin{explain}
सातत्य परिकल्पना वह कथन है कि सातत्य का आकार पहली !!{nonenumerable}
कार्डिनैलिटी है, अर्थात् $\Pow{\Nat}$ का आकार~$\aleph_1$ है।
\end{explain}

\begin{prop}
सातत्य परिकल्पना तभी और केवल तभी सत्य है जब \[\fn{CH} \ident
\lforall[X][(\fn{Aleph}_1(X) \liff \fn{Cont}(X))]\] वैध हो।
\end{prop}

ध्यान दें कि यह सत्य नहीं है कि $\lnot \fn{CH}$ तभी और केवल तभी वैध है जब
सातत्य परिकल्पना असत्य हो। किसी !!a{enumerable} प्रांत में न तो
आकार~$\aleph_1$ के उपसमुच्चय होते हैं, न सातत्य के आकार के उपसमुच्चय;
इसलिए किसी !!a{enumerable} प्रांत में $\fn{CH}$ सर्वदा सत्य है। फिर भी, हम
ऐसा भिन्न वाक्य दे सकते हैं जो तभी और केवल तभी वैध है जब सातत्य परिकल्पना
असत्य हो:

\begin{prop}
सातत्य परिकल्पना तभी और केवल तभी असत्य है जब \[\fn{NCH} \ident
\lforall[X][(\fn{Cont}(X) \lif \lexists[Y][(Y \subseteq X \land \lnot
    \fn{Count}(Y) \land \lnot \cardeq{X}{Y})])]\] वैध हो।
\end{prop}

\end{document}
